\documentclass[12pt,a4paper]{article}
\usepackage[UTF8]{ctex}
\usepackage{amsmath,amssymb,geometry,xcolor,listings,tcolorbox}
\geometry{margin=2cm}
\lstset{language=C,basicstyle=\ttfamily\small,frame=single,numbers=left,breaklines=true}
\newtcolorbox{keybox}[1][]{colback=yellow!5!white,colframe=yellow!75!black,title=考点,#1}

\title{FDS 复习 — 并查集 + 图论}
\author{6月29日 冲刺}
\date{}

\begin{document}
\maketitle

\part{并查集 (Union-Find / Disjoint Set)}

\section{核心问题}

\begin{quote}
给定 N 个元素，支持两种操作：
\begin{itemize}
\item \textbf{Find(x):} x 属于哪个集合？（找代表元）
\item \textbf{Union(x,y):} 把 x 和 y 所在的两个集合合并
\end{itemize}
\end{quote}

\textbf{应用：} 动态等价关系、Kruskal 最小生成树、连通分量检测。

\section{数据结构}

\begin{lstlisting}
int S[MAX];   // S[root] = -size  (负数表示集合大小)
              // S[child] = parent index (正数表示父节点)
\end{lstlisting}

\subsection{初始化}

\begin{lstlisting}
void init(int N) {
    for (int i = 0; i < N; i++) S[i] = -1;  // 每个元素自成集合
}
\end{lstlisting}

\subsection{Find — 路径压缩}

\begin{lstlisting}
int find(int x) {
    if (S[x] < 0) return x;          // x 是根
    return S[x] = find(S[x]);        // 爹变祖宗！压缩路径
}
\end{lstlisting}

\textbf{效果：} 查找路径上所有节点直接指向根，后续查找几乎 $O(1)$。

\subsection{Union — 按大小合并}

\begin{lstlisting}
void unionSet(int a, int b) {
    a = find(a); b = find(b);
    if (a == b) return;
    if (S[a] > S[b]) { int t = a; a = b; b = t; }  // a 更大
    S[a] += S[b];    // 更新大小 (负得更多)
    S[b] = a;        // b 挂到 a
}
\end{lstlisting}

\textbf{思路：} 总是把小的树挂到大的树下，树高不超过 $\lfloor \log_2 N\rfloor + 1$。

\subsection{复杂度}

\begin{keybox}
\textbf{Union-by-Size + 路径压缩：}
$$T(M, N) = O(M \cdot \alpha(M, N))$$
$\alpha$ 是反 Ackermann 函数，$\alpha \leq 4$ 在实际数据范围内。

\textbf{几乎就是常数时间！}
\end{keybox}

\begin{itemize}
\item 仅 Union-by-Size：$T = O(N + M \log N)$
\item 朴素 Union（不压缩、不按大小）：$\Theta(N^2)$ 最坏
\end{itemize}

\section{四个版本对比}

\begin{tabular}{|c|c|c|}
\hline
\textbf{优化} & \textbf{Find} & \textbf{Union} \\
\hline
0. 朴素 & 一路走到根 & 随便挂 \\
1. Union-by-Size & 朴素 Find & 小挂大 \\
2. 路径压缩 & 爹变祖宗 & 随便挂 \\
3. 两者都有 & 路径压缩 & 小挂大 \\
\hline
\end{tabular}

版本 3 是标准写法。

\begin{keybox}
\textbf{路径压缩不能和 Union-by-Height 一起用！} 压缩会改变高度，所以改成 Union-by-Rank（秩 = 估计高度）。
\end{keybox}

\part{图论 (Graph)}

\section{基本概念}

\begin{itemize}
\item $G = (V, E)$，$|V|$ 顶点数，$|E|$ 边数
\item \textbf{无向图:} $(v_i, v_j) = (v_j, v_i)$
\item \textbf{有向图:} $\langle v_i, v_j \rangle \neq \langle v_j, v_i \rangle$
\item \textbf{DAG:} 有向无环图
\item \textbf{握手定理:} $\sum degree(v) = 2|E|$
\end{itemize}

\section{图的表示}

\subsection{邻接矩阵}
\begin{itemize}
\item $O(|V|^2)$ 空间
\item 适合稠密图
\item 判断两顶点是否相邻：$O(1)$
\end{itemize}

\subsection{邻接表}
\begin{itemize}
\item $O(|V| + |E|)$ 空间
\item 适合稀疏图
\item 遍历所有边：$O(|V| + |E|)$
\end{itemize}

\begin{keybox}
\textbf{考试判断：} 稠密图用矩阵，稀疏图用表。
\end{keybox}

\section{拓扑排序}

\textbf{目标：} 给 DAG 的顶点排一个线性顺序，满足所有边的方向（前驱在前）。

\begin{enumerate}
\item 计算所有顶点入度
\item 入度 = 0 的进队列
\item 出队 → 所有邻接入度 -1 → 减到 0 的进队
\item 重复直到队列空
\item 如果输出的顶点数 < |V| → 图中有环！
\end{enumerate}

\textbf{时间复杂度：} $O(|V| + |E|)$

\begin{lstlisting}
void topsort(Graph G) {
    Queue Q;
    for each V
        if (Indegree[V] == 0) Enqueue(V, Q);
    int cnt = 0;
    while (!IsEmpty(Q)) {
        V = Dequeue(Q);
        TopNum[V] = ++cnt;
        for each W adjacent to V
            if (--Indegree[W] == 0) Enqueue(W, Q);
    }
    if (cnt < NumVertex) Error("Graph has a cycle");
}
\end{lstlisting}

\section{最短路径}

\subsection{无权最短路径 — BFS}

BFS 就是无权图的最短路径算法。$O(|V| + |E|)$。

\subsection{Dijkstra — 带权最短路径}

\textbf{思想：} 贪心。每次选 Dist 最小的未处理顶点，松弛所有邻接点。

\textbf{条件：不能有负权边！}

\begin{lstlisting}
void dijkstra(int graph[][MAX], int S, int N) {
    int known[MAX] = {0};
    int dist[MAX];
    for (int i = 0; i < N; i++) dist[i] = INF;
    dist[S] = 0;
    for (;;) {
        // 找未处理中距离最小的
        int V = -1, minD = INF;
        for (int i = 0; i < N; i++)
            if (!known[i] && dist[i] < minD)
                { minD = dist[i]; V = i; }
        if (V == -1) break;
        known[V] = 1;
        // 松弛所有邻接点
        for (int W = 0; W < N; W++)
            if (graph[V][W] && !known[W])
                if (dist[V] + graph[V][W] < dist[W])
                    dist[W] = dist[V] + graph[V][W];
    }
}
\end{lstlisting}

\textbf{复杂度：} 朴素 $O(|V|^2)$，用堆优化 $O(|E| \log |V|)$。

\subsection{负权边 — Bellman-Ford}

每轮松弛所有边，共 $|V|-1$ 轮。$O(|V| \cdot |E|)$。

\textbf{第 |V| 轮如果还能松弛 → 有负权环！}

\subsection{全源最短路径 — Floyd-Warshall}

\begin{lstlisting}
for (int k = 0; k < N; k++)
    for (int i = 0; i < N; i++)
        for (int j = 0; j < N; j++)
            if (dist[i][k] + dist[k][j] < dist[i][j])
                dist[i][j] = dist[i][k] + dist[k][j];
\end{lstlisting}

$O(|V|^3)$。

\section{最小生成树 (MST)}

\subsection{Prim}

\textbf{思想：} 类似 Dijkstra，但 Dist 是从当前树到各顶点的最短距离。

\textbf{复杂度：} 朴素 $O(|V|^2)$，堆优化 $O(|E| \log |V|)$。

\subsection{Kruskal}

\textbf{思想：} 排序所有边 + Union-Find 判环。

\begin{lstlisting}
void kruskal(Graph G) {
    T = {};
    sort(edges by weight);
    while (T has < |V|-1 edges) {
        pick min edge (u,v);
        if (find(u) != find(v)) {  // 不同集合 → 不成环
            add to T;
            union(u, v);
        }
    }
}
\end{lstlisting}

\textbf{复杂度：} $O(|E| \log |E|) = O(|E| \log |V|)$。

\begin{keybox}
\textbf{Prim vs Kruskal:}
\begin{itemize}
\item Prim：类似 Dijkstra，适合\textbf{稠密图}
\item Kruskal：排序边 + 并查集，适合\textbf{稀疏图}
\end{itemize}
\end{keybox}

\section{AOE 网络与关键路径}

\begin{itemize}
\item \textbf{EC[j]} = 活动 j 的最早完成时间 = $\max_{i \to j}(EC[i] + C_{ij})$
\item \textbf{LC[i]} = 活动 i 的最晚完成时间 = $\min_{i \to j}(LC[j] - C_{ij})$
\item \textbf{Slack Time} = LC - EC，为 0 的边在\textbf{关键路径}上
\end{itemize}

\section{网络流 — Ford-Fulkerson}

\textbf{核心技巧：残差图 + 反向边（撤销机制）。}

\begin{enumerate}
\item 在残差图中找一条增广路
\item 增加正向边流量，减少反向边（等于增加残差）
\item 重复直到不存在增广路
\end{enumerate}

\section{DFS 应用}

\subsection{双连通分量与关节点}

\begin{itemize}
\item \texttt{Num[v]} = DFS 编号
\item \texttt{Low[v]} = 通过回边能到达的最小 Num
\item \textbf{关节点条件：} $Low[child] \ge Num[v]$，且 v 不是根 / v 是根且有 $\ge 2$ 个孩子
\end{itemize}

\subsection{欧拉回路}

\begin{itemize}
\item 无向图有欧拉回路 $\iff$ 所有顶点度数为偶数
\item 无向图有欧拉路径 $\iff$ 恰好两个顶点度数为奇数
\item 有向图有欧拉回路 $\iff$ 每个顶点入度 = 出度
\end{itemize}

\end{document}
